AI creative generators are stochastic: the same pipeline, given the same inputs, produces
different outputs across independent runs. Practitioners deploying these systems lack principled
guidance on which design decisions will be stable across runs---and can therefore be trusted
on a single execution---and which will vary, requiring multiple runs and human selection
to control quality.

We address this problem through a \emph{same-input divergence} study: two fully independent,
clean-room runs of a complete AI puzzle hunt generation pipeline from identical inputs---the
same Age of Empires~II world data, scope document, and toolkit specification---compared at
five levels of abstraction (structural, pool, answer-word assignment, extraction mechanism,
and micro-design). An AI expert panel evaluated both outputs on a 7-dimension, 45-point
weighted rubric using identical reviewer configurations.

The results reveal a structured, not random, pattern of divergence. Every decision explicitly
specified in the creative brief---scale (5 feeders + 1 meta), narrator voice, audience,
format, and solve-time target---is identical across both runs. Every decision left to the
pipeline diverges substantially: round count (1 vs.~3), meta mechanism family (crossword
grid vs.~vowel-count index followed by anagram resolution), and answer words
(Jaccard~$= 0.00$ across all six words including the meta). We identify a distinct
phenomenon, \emph{domain semantic gravity}: despite zero lexical overlap, 3 of 5 feeder
slots yield domain-adjacent word pairs (e.g., \textsc{onager}/\textsc{siege};
\textsc{patrol}/\textsc{march}), revealing that the domain exerts a semantic basin of
attraction without constraining specific word choice. Panel quality scores are stable at
the pass/fail level (all authored puzzles pass in both runs; normalized aggregate scores
within 3 percentage points) while showing a systematic $+1.2$-point uplift in Elegance
and Reading Reward in the second run, consistent with a genuine mechanism-level improvement.

We introduce same-input divergence as a generalizable experimental paradigm for
characterizing AI pipeline reliability, establish \emph{specification drives convergence}
as the primary practical principle---every under-specified decision is a variable
decision---and derive actionable deployment guidance for practitioners operating AI creative
pipelines at scale.
